% Generated draft: requires mathematical and source review.
\documentclass[11pt]{article}
\usepackage[margin=25mm]{geometry}
\usepackage{fontspec}
\setmainfont{Latin Modern Roman}
\usepackage{amsmath,amssymb,mathtools,unicode-math}
\setmathfont{Latin Modern Math}
\usepackage{enumitem,needspace,xurl,hyperref}
\setlength{\parindent}{0pt}
\setlength{\parskip}{5pt}
\setlength{\emergencystretch}{3em}
\newcommand{\N}{\mathbb{N}}
\newcommand{\Z}{\mathbb{Z}}
\newcommand{\Q}{\mathbb{Q}}
\newcommand{\R}{\mathbb{R}}
\newcommand{\C}{\mathbb{C}}
\newcommand{\F}{\mathbb{F}}
\newcommand{\A}{\mathbb{A}}
\newcommand{\PP}{\mathbb{P}}
\newcommand{\E}{\mathbb{E}}
\newcommand{\eps}{\varepsilon}
\newcommand{\dd}{\,\mathrm{d}}
\begin{document}
\section*{UnsolvedMath ID: 30000518}
\textbf{Source problem code:} catalogue entry
\subsection{Definitions and mathematical statement}
Let \((\mathcal E_k,\mathcal R_k)\) be the source's energy/evolution-functional pairs, or \((\mathcal E_k,\mathcal K_k)\) in its Hamiltonian formulation, with intended independent \(\Gamma\)-limits \((\mathcal E,\mathcal R)\) or \((\mathcal E,\mathcal K)\). Let \(q_k\) be corresponding solutions on a time interval \(I\), and suppose \(q_k(t)\to q(t)\) for every \(t\in I\) in the intended configuration topology.

For reference, sequential \(\Gamma\)-convergence of \(F_k\) to \(F\) on a topological domain requires
\[
F(y)\leq\liminf_k F_k(y_k)\quad\text{for every }y_k\to y,
\qquad
\forall y\ \exists y_k\to y:\ \limsup_kF_k(y_k)\leq F(y).
\]
For time-dependent functionals the record does not say whether these statements hold at each fixed time, uniformly in time, or for integral functionals on paths. It also does not identify the domains of the second functionals, which may involve velocity or other evolution data rather than just configuration.

The research question is to identify joint convergence and compatibility hypotheses under which the limiting path is a solution of the evolution determined by \((\mathcal E,\mathcal R)\), respectively \((\mathcal E,\mathcal K)\). The two formulations are distinct and both are retained. Independent \(\Gamma\)-convergence of the components and pointwise path convergence are given data, not claimed sufficient. Separate recovery sequences need not be compatible with each other or with the solution path, and convergence of static variational values alone does not specify convergence of the evolutionary relations.

A precise theorem needs the configuration and path spaces, time dependence, definitions of \(\mathcal R\) and \(\mathcal K\), solution concept, initial conditions and energy preparation, and the proposed joint liminf/recovery or dynamical compatibility conditions. No particular gradient-flow equation, Hamiltonian equation, subdifferential, dissipation balance, or kinetic interpretation of \(\mathcal K\) is imposed without those definitions.

The supplied assessment checked 2026-08-21 describes the fragment as lacking enough hypotheses to verify a determinate resolution. The cleaned statement names the limiting pairs, but it still gives a programme for sufficient joint conditions rather than one specified general convergence theorem.
\subsection{Short English statement}
Find joint variational-convergence conditions that let limits of solutions solve the limiting evolution, for both dissipation-type and Hamiltonian pairs. Separate \(\Gamma\)-limits and pointwise convergence are not asserted sufficient; the dynamics and compatibility conditions remain unspecified.
\subsection{Sources}
\begin{itemize}
\item Local UnsolvedMath registry, permanent ID 30000518, catalogue entry, \emph{Joint \(\Gamma\)-Convergence for Evolutionary Variational Problems}.
\item \emph{Applications of Asymptotic Analysis} (2007), Oberwolfach Reports, pp. 1663--1730, DOI \url{https://doi.org/10.4171/owr/2006/28}.
\item Supplied source-verification claim and literature triage checked 2026-08-21 cite only that report; no complete dynamical definitions, joint convergence criterion, or separate theorem is supplied.
\end{itemize}
\end{document}
